% Created 2024-01-05 vie 20:26
% Intended LaTeX compiler: pdflatex
\documentclass[11pt]{article}
\usepackage[utf8]{inputenc}
\usepackage[T1]{fontenc}
\usepackage{graphicx}
\usepackage{longtable}
\usepackage{wrapfig}
\usepackage{rotating}
\usepackage[normalem]{ulem}
\usepackage{amsmath}
\usepackage{amssymb}
\usepackage{capt-of}
\usepackage{hyperref}
\usepackage{minted}

\usepackage{parskip}
\usepackage[utf8]{inputenc} %% For unicode chars
\usepackage{placeins}
\usepackage[margin=2.50cm]{geometry}
\usepackage[T1]{fontenc}
\usepackage{mathpazo}
\linespread{1.05}
\usepackage[scaled]{helvet}
\usepackage{courier}
\hypersetup{colorlinks=true,linkcolor=blue}
\RequirePackage{fancyvrb}
\usepackage{mdframed}
\BeforeBeginEnvironment{minted}{\begin{mdframed}}
\AfterEndEnvironment{minted}{\end{mdframed}}
\setcounter{secnumdepth}{3}
\date{}
\title{}
\hypersetup{
 pdfauthor={},
 pdftitle={},
 pdfkeywords={},
 pdfsubject={},
 pdfcreator={Emacs 28.1 (Org mode 9.6.12)}, 
 pdflang={Spanish}}
\begin{document}

\setcounter{tocdepth}{3}
\tableofcontents

\setlength\parindent{10pt}

\section{Chapter 4: Loops in PHP}
\label{sec:org7c729b7}
Loops are a fundamental part of programming in PHP, allowing you to
execute a block of code repeatedly under certain conditions. This
ability to iterate over data or execute code multiple times is crucial
in many programming scenarios. In this chapter, we will cover the four
primary types of loops in PHP: \texttt{for}, \texttt{foreach}, \texttt{while}, and
\texttt{do-while}, along with guidance on when to use each type for optimal
results.

\subsection{For Loop}
\label{sec:org445c697}
The \texttt{for} loop is ideal when you know in advance how many times you want
to execute a statement or a block of statements.

\subsubsection{Syntax}
\label{sec:orgd817606}
\begin{minted}[]{php}
for (initialization; condition; increment) {
    // Code to be executed
}
\end{minted}

\subsubsection{Example}
\label{sec:orgefe7679}
For example, to print numbers from 1 to 10:

\begin{minted}[]{php}
for ($i = 0; $i < 10; $i++) {
    echo $i . " ";
}
\end{minted}
\begin{minted}[]{php}
for ($i = 1; $i <= 10; $i++) {
    echo $i . "\n";
}
\end{minted}

\subsubsection{When to Use}
\label{sec:org479e310}
\begin{itemize}
\item \textbf{Count-Controlled Iteration}: Ideal for looping a specific number of
times.
\item \textbf{Array Indexing}: Useful for iterating over array elements when the
index is needed.
\end{itemize}

\subsection{Foreach Loop}
\label{sec:orgc46f589}
The =foreach=` loop is specifically designed for iterating over arrays
or objects, making it simpler to work with elements without worrying
about indexes or keys.

\subsubsection{Syntax}
\label{sec:org952f03f}
\begin{minted}[]{php}
foreach ($array as $value) {
    // Code to be executed
}
\end{minted}

Or, to access both the key and the value:

\begin{minted}[]{php}
foreach ($array as $key => $value) {
    // Code to be executed
}
\end{minted}

\subsubsection{Example}
\label{sec:org1e4a04e}
To print all elements of an array>
\begin{minted}[]{php}
$colors = ["red", "green", "blue"];

foreach ($colors as $color) {
    echo $color . " ";
}
\end{minted}


To print all values of an associative array:

\begin{minted}[]{php}
$colors = ["red" => "#FF0000", "green" => "#00FF00", "blue" => "#0000FF"];

foreach ($colors as $color) {
    echo $color . "\n";
}
\end{minted}

\subsubsection{When to Use}
\label{sec:org8db14a5}
\begin{itemize}
\item \textbf{Array Iteration}: for iterating through arrays or objects where the
order and number of elements are unknown or irrelevant.
\item \textbf{Readable Code}: Offers clearer and more readable code when dealing
with array or objects elements.
\end{itemize}

\subsection{While Loop}
\label{sec:org25c57a9}
The \texttt{while} loop is used when you want to execute a block of code as
long as a specified condition is true.

\subsubsection{Syntax}
\label{sec:org3869f47}
\begin{minted}[]{php}
while (condition) {
    // Code to be executed
}
\end{minted}

\subsubsection{Example}
\label{sec:org17e8aa8}
\begin{minted}[]{php}
$i = 0;
while ($i < 10) {
    echo $i . " ";
    $i++;
}
\end{minted}

\subsubsection{When to Use}
\label{sec:org721c488}
\begin{itemize}
\item \textbf{Unknown Iterations}: Suitable when the number of iterations is not
known in advance.
\item \textbf{Conditional Control}: Useful when the continuation condition might be
affected by processes within the loop.
\end{itemize}

\subsection{Do-While Loop}
\label{sec:org22f33f4}
The \texttt{do-while} loop is a variant of the while loop. This loop will
execute the code block once before checking if the condition is true,
then it will repeat the loop as long as the condition is true.

\subsubsection{Syntax}
\label{sec:orgfa2210b}
\begin{minted}[]{php}
do {
    // Code to be executed
} while (condition);
\end{minted}

\subsubsection{Example}
\label{sec:orgd497704}
\begin{minted}[]{php}
$i = 0;
do {
    echo $i . " ";
    $i++;
} while ($i < 10);
\end{minted}

\subsubsection{When to Use}
\label{sec:orgb797602}
\begin{itemize}
\item \textbf{Guaranteed Execution}: Ensures that the loop body is executed at
least once.
\item \textbf{Post-Condition Looping}: Ideal when the condition for continuation is
determined at the end of each iteration.
\end{itemize}

\subsubsection{Conclusion}
\label{sec:orga573b74}
Understanding the different types of loops and their appropriate use
cases is crucial for efficient PHP programming. While \texttt{for} loops are
best for fixed iterations, \texttt{foreach} is optimal for array traversal
without needing to know the array size. \texttt{while} loops suit scenarios
where the number of iterations is not predetermined, and \texttt{do-while}
loops guarantee at least one iteration, regardless of the condition. As
you continue to develop your PHP skills, experiment with these loops to
find the best fit for various coding challenges.

\section{Loops alternative syntax}
\label{sec:org455c5bf}
In PHP, the alternative syntax for loops provides a way to write loops
that is often more readable, especially when embedding PHP within HTML.
This syntax is available in PHP 8 and earlier versions. The alternative
syntax is particularly useful in templating where PHP is interspersed
with HTML.

\subsection{Alternative Syntax for \texttt{for} Loop}
\label{sec:org215a965}
Instead of using curly braces \texttt{\{\}}, you use \texttt{:} to start and \texttt{endfor;}
to end the loop.

\textbf{Standard Syntax:}

\begin{minted}[]{php}
for ($i = 0; $i < 10; $i++) {
    // Code to execute
}
\end{minted}

\textbf{Alternative Syntax:}

\begin{minted}[]{php}
for ($i = 0; $i < 10; $i++):
    // Code to execute
endfor;
\end{minted}

\subsection{Alternative Syntax for \texttt{foreach} Loop}
\label{sec:org9c91092}
Similar to the \texttt{for} loop, the \texttt{foreach} loop can use \texttt{endforeach;} in
place of curly braces.

\textbf{Standard Syntax:}

\begin{minted}[]{php}
foreach ($array as $element) {
    // Code to execute
}
\end{minted}

\textbf{Alternative Syntax:}

\begin{minted}[]{php}
foreach ($array as $element):
    // Code to execute
endforeach;
\end{minted}

\subsection{Alternative Syntax for \texttt{while} Loop}
\label{sec:org76e70f1}
The \texttt{while} loop can use \texttt{endwhile;} at the end instead of curly braces.

\textbf{Standard Syntax:}

\begin{minted}[]{php}
while ($condition) {
    // Code to execute
}
\end{minted}

\textbf{Alternative Syntax:}

\begin{minted}[]{php}
while ($condition):
    // Code to execute
endwhile;
\end{minted}

\subsection{Alternative Syntax for \texttt{do-while} Loop}
\label{sec:org6833352}
The \texttt{do-while} loop can also be written using alternative syntax, though
it's less common.

\textbf{Standard Syntax:}

\begin{minted}[]{php}
do {
    // Code to execute
} while ($condition);
\end{minted}

\textbf{Alternative Syntax:}

\begin{minted}[]{php}
do:
    // Code to execute
while ($condition);
\end{minted}

\subsection{Example in an HTML Template}
\label{sec:org56b8d1b}
Here's how you might use the alternative syntax in an HTML template:

\begin{minted}[]{php}
<ul>
<?php foreach ($items as $item): ?>
    <li><?php echo htmlspecialchars($item); ?></li>
<?php endforeach; ?>
</ul>
\end{minted}

This example demonstrates how the alternative syntax can make PHP code
embedded in HTML more readable.

\subsection{Conclusion}
\label{sec:org112aae7}
The alternative syntax for loops in PHP doesn't change the functionality
of the loops but offers a way to write code that can be easier to read
and maintain, especially in the context of PHP/HTML templates. This
syntax is a matter of preference and style, and its usage often depends
on the coding standards of a particular project or team.

\section{Break and Continue in PHP}
\label{sec:org1b1f0c3}
In PHP, \texttt{break} and \texttt{continue} are important control structures used
within loops (\texttt{for}, \texttt{foreach}, \texttt{while}, \texttt{do-while}) to alter the normal
flow of the loop. Understanding how to use them correctly can make your
loops more efficient and powerful.

\subsection{The \texttt{break} Statement}
\label{sec:orgf1f1089}
\texttt{break} is used to terminate the execution of a loop prematurely. When
\texttt{break} is encountered inside a loop, the loop is immediately stopped,
and the program control resumes at the next statement following the
loop.

\subsubsection{Use Case of \texttt{break}:}
\label{sec:org5b9843f}
\begin{itemize}
\item To exit a loop when a certain condition is met.
\item Commonly used in switch-case statements to terminate a case.
\end{itemize}

\subsubsection{Example:}
\label{sec:orge69709b}
\begin{verbatim}
for ($i = 0; $i < 10; $i++) {
    if ($i == 5) {
        break; // Exit the loop when $i equals 5
    }
    echo $i . " ";
}
// Output: 0 1 2 3 4
\end{verbatim}

\subsection{The \texttt{continue} Statement}
\label{sec:org1c6713a}
\texttt{continue} is used to skip the rest of the current loop iteration and
proceed with the next iteration. Unlike \texttt{break}, \texttt{continue} does not
terminate the loop but just skips the remaining statements in the
current iteration and continues with the next iteration.

\subsubsection{Use Case of \texttt{continue}:}
\label{sec:org4d1b647}
\begin{itemize}
\item To skip certain iterations based on a condition.
\item Useful for avoiding nested if-else statements inside loops.
\end{itemize}

\subsubsection{Example:}
\label{sec:orgffe7985}
\begin{verbatim}
for ($i = 0; $i < 10; $i++) {
    if ($i % 2 == 0) {
        continue; // Skip the rest of the loop for even numbers
    }
    echo $i . " ";
}
// Output: 1 3 5 7 9
\end{verbatim}

\subsection{Note on Using \texttt{break} and \texttt{continue}}
\label{sec:orgd002053}
\begin{itemize}
\item \texttt{break} and \texttt{continue} only affect the innermost loop that encloses
them.
\item In nested loops, you can specify a numeric argument to indicate how
many levels of enclosing loops should be broken out of or continued.
For example, \texttt{break 2;} will break out of two levels of loops.
\end{itemize}

\subsection{Conclusion}
\label{sec:orgbae73b6}
\texttt{break} and \texttt{continue} are powerful tools for controlling the execution
flow inside loops. They can make your code more readable and efficient
by preventing unnecessary iterations or exiting a loop as soon as a
certain condition is met. However, they should be used judiciously to
keep the logic of your loops clear and understandable.
\end{document}
